The device model of Sec.~\ref{sec:model} fixed the per-sensor chain from a sensor share $E_s$ to the peak tunneling rate $\Gamma_{\rm in}^{\rm pk}$ and the saturation ceiling $\Gamma_{\rm max}=1/\tau_d=25$~kHz [Eqs.~\eqref{eq:chain}--\eqref{eq:censor}], but deferred two ingredients to this section: how the deposited energy $E_{\rm dep}$ is partitioned among the $\sim$10,300 sensors, and how a reconstructed energy $E_{\rm rec}$ is estimated from the resulting censored counts. We specify both here and use them to build the Monte-Carlo response matrix $R(E_{\rm rec}\,|\,E_{\rm dep})$ that Sec.~\ref{sec:results} folds through the deposit spectra. This is the methodological core of the pipeline, and also its principal source of model risk: the saturated-regime response has no published benchmark, and the crossover between the linear and saturated regimes rides on two low-confidence sharing parameters. We are explicit about both below.

\emph{Sensor-sharing partition.} A single deposit does not illuminate the $N_{\rm sens}\simeq10{,}300$ sensors uniformly. We adopt a localized-plus-diffuse partition: a prompt fraction $f_{\rm prompt}$ of $E_{\rm dep}$ is absorbed near the impact point and spread over a spot of $\sim$$\pi r^2$ sensors, while the remaining $(1-f_{\rm prompt})$ thermalizes diffusely across the whole array. The on-spot and off-spot per-sensor shares are then
\begin{align}
E_s^{\rm on}  &= \frac{f_{\rm prompt}\,E_{\rm dep}}{\pi r^2} + \frac{(1-f_{\rm prompt})\,E_{\rm dep}}{N_{\rm sens}}, \nonumber\\
E_s^{\rm off} &= \frac{(1-f_{\rm prompt})\,E_{\rm dep}}{N_{\rm sens}},
\label{eq:sharing}
\end{align}
where $r$ is the spot radius in sensor-pitch units. Both sides carry units of energy. The two parameters $f_{\rm prompt}$ (default $0.3$, range $0.1$--$0.5$) and $r$ (default $2$, range $1$--$5$) have no direct thin-wafer QPD measurement; we treat them as exposed, scanned inputs rather than derived values, and we carry the equal-split limit $f_{\rm prompt}\to0$ (uniform sharing) as a required sanity bound. Because $E_s^{\rm on}$ is the largest per-sensor share, the on-spot sensors are the first to saturate, and $f_{\rm prompt}$ and $r$ set the deposit energy at which that happens.

\emph{Count-integral estimator.} On each sensor the trapped-quasiparticle population $N_{\rm qp}(E_s)$ [Eq.~\eqref{eq:chain}] drives a two-exponential tunneling burst. The quantity that controls both the burst statistics and the saturation is the expected number of \emph{tunneling events}, $n_{\rm ev}=\int\Gamma_{\rm in}(t)\,dt=K\tau_{\rm qp}N_{\rm qp}/V_{\rm tr}$, where $\tau_{\rm qp}$ is the trapped-quasiparticle lifetime; this is distinct from, and much smaller than, the trapped count $N_{\rm qp}$ itself (by $\sim$$0.03$ for Al, $\sim$$0.008$ for Hf), and it is $n_{\rm ev}$, not $N_{\rm qp}$, that we feed as the Poisson mean to the exponentially-modified-Gaussian (EMG) tunneling-burst template of Ref.~\cite{ramanathan2026}. Conflating the two would mis-scale the entire saturation. The template realizes the event train, which we censor with the dead-time model $m(\Gamma_{\rm in})$ [Eq.~\eqref{eq:censor}] to obtain a per-sensor observed count $N_{{\rm obs},i}$; summing over the array gives the reconstructed energy
\begin{equation}
\begin{aligned}
E_{\rm rec}&=C\,N_{\rm obs},\qquad
N_{\rm obs}=\sum_{i=1}^{N_{\rm sens}} N_{{\rm obs},i},\\
N_{{\rm obs},i}&=\int m\big(\Gamma_{{\rm in},i}(t)\big)\,dt .
\end{aligned}
\label{eq:erec}
\end{equation}
The count-to-energy constant $C$ (units eV per event) is a \emph{single global} constant per design, fixed once on a deeply unsaturated deposit by requiring the linear slope $dE_{\rm rec}/dE_{\rm dep}=\varepsilon=0.5$. It is not tuned per energy bin. Consequently, the recovery of $E_{\rm rec}\approx0.5\,E_{\rm dep}$ at low energy is a calibration-consistency check, not an independent validation; the genuinely predictive content of Eq.~\eqref{eq:erec} is the saturated-regime behavior, to which we turn next.

\emph{Saturation and the crossover band.} When the on-spot peak rate reaches the ceiling, $\Gamma_{\rm in}^{\rm pk}=\Gamma_{\rm max}$ [Eq.~\eqref{eq:sat}], the censored count $N_{{\rm obs},i}$ stops tracking $n_{\rm ev}$: under non-paralyzable censoring it \emph{plateaus} at $\sim$$1/\tau_d$ per sensor, so the summed count in Eq.~\eqref{eq:erec} plateaus and $E_{\rm rec}$ saturates. This plateau \emph{is} the modeled bandwidth saturation; at no point do we linearly extrapolate a rate to an energy. The retained paralyzable variant instead rolls the count over (the observed rate peaks at $\Gamma=\Gamma_{\rm max}$ and declines), which we carry only as a sensitivity. Writing $E_{\rm onset}$ for the per-sensor share at which $\Gamma_{\rm in}^{\rm pk}=\Gamma_{\rm max}$ ($\sim$$1.27$~eV for Ta$\to$Al, $\sim$$0.77$~eV for Al$\to$Hf, from Eqs.~\eqref{eq:chain}--\eqref{eq:sat}), the crossover deposit energy at which the on-spot sensor first saturates follows from inverting Eq.~\eqref{eq:sharing}:
\begin{equation}
E_{\rm dep}^{\rm cross}=\frac{E_{\rm onset}}
{\dfrac{f_{\rm prompt}}{\pi r^2}+\dfrac{1-f_{\rm prompt}}{N_{\rm sens}}} .
\label{eq:crossover}
\end{equation}
Because $E_{\rm dep}^{\rm cross}$ inherits the $f_{\rm prompt}/r$ uncertainty, we report it as a band rather than a single number. At the default point it is $\sim$52.9~eV (Ta$\to$Al) and $\sim$32.1~eV (Al$\to$Hf); over the full $f_{\rm prompt}\in[0.1,0.5]$, $r\in[1,5]$ scan it ranges up to $\sim$13~keV and $\sim$7.9~keV in the equal-split limit, with a whole-array plateau (every sensor saturated) at $\sim$18.6~keV and $\sim$11.3~keV. The Al$\to$Hf design saturates before Ta$\to$Al at every point of the scan, consistent with its lower per-sensor onset. These are the SIMU-01 design numbers forward-referenced in Sec.~\ref{sec:model}; the spread of roughly three orders of magnitude is $f_{\rm prompt}$-dominated and is an honest reflection of the sharing uncertainty, not a computed precision.

\emph{Response matrix.} We assemble the bandwidth-limited transfer function
\begin{equation}
\begin{aligned}
R(E_{\rm rec}\,|\,E_{\rm dep})&=\frac{dP}{dE_{\rm rec}}\bigg|_{E_{\rm dep}},\\
\int R(E_{\rm rec}\,|\,E_{\rm dep})\,dE_{\rm rec}&=1,
\end{aligned}
\label{eq:response}
\end{equation}
the normalized conditional distribution of reconstructed energy for a fixed deposit, per design and per censoring variant. Each $E_{\rm dep}$ column is built by Monte-Carlo, histogramming $N_s=5000$ forward realizations of Eq.~\eqref{eq:erec} over a grid of $584$ deposit energies sampled uniformly in $\log E_{\rm dep}$ from $10.14$~eV to $197$~MeV (80 bins per decade), spanning the sub-keV CEvNS floor through the deep-saturated cosmic-muon tail. The $\sim$$\pi r^2$ on-spot sensors are realized individually with the EMG template; the $\sim$10,287 identical off-spot sensors are drawn as one aggregate ensemble per realization rather than looped. Columns normalize to unity to $10^{-12}$, and the well-populated cells converge to $\lesssim$1.8\% at $N_s=5000$ with clean $1/\sqrt{N_s}$ scaling. Figure~\ref{fig:response} shows the resulting median $E_{\rm rec}(E_{\rm dep})$ with its 16--84\% band for both designs and both variants. At the $197$~MeV muon endpoint the non-paralyzable response plateaus at $\sim$34.8~keV (Ta$\to$Al) and $\sim$26.8~keV (Al$\to$Hf)---some three to four orders of magnitude below the linear $0.5\,E_{\rm dep}$ line ($\approx$98.5~MeV)---while the paralyzable variant rolls over to $\sim$3.3~keV and $\sim$2.6~keV. No column ramps linearly into the MeV range.

\emph{Limiting-case validation and its limits.} The estimator is validated only at its two limits (VALD-04): $E_{\rm rec}\approx0.5\,E_{\rm dep}$ below the crossover, and a plateau once the peak tunneling rate exceeds $25$~kHz, driven by the censoring model rather than by any numerical artifact (with censoring disabled, the count integral returns to exactly $0.5\,E_{\rm dep}$, so the plateau demonstrably has teeth). We state plainly that \emph{there is no literature anchor for the shape of the response between these two limits}: the saturated-regime curve is a model prediction, not a benchmarked result. The tens-of-keV muon pile-up that this section predicts is therefore an instrumental artifact of the modeled $25$~kHz ceiling---not a physical spectral feature---and should be read as ``the model says,'' not ``the detector measures.'' With this caveat in hand, Eq.~\eqref{eq:response} is the object we fold the deposited-energy spectra through in Sec.~\ref{sec:results}.
